% !TeX root = ../../Book3.tex
\section{同调形式}
\label{b3:sec:1003}

同调用可加的圈代替需要顺序的路径类。在平面中，零同调可以用区域外所有点的指数为零来判断。格网分割将这个条件转换成小方形边界的相消，从而给全局积分公式。

\subsection{圈}
\label{b3:subsec:100301}

一个圈 \(\Gamma\) 是有限个分段光滑闭曲线的整数线性组合 \(\sum_j n_j\gamma_j\)。积分、指数均按系数线性定义。若先给有限条有向弧并要求每点的端点系数总和为0，可把这些弧依次拼成闭路，得到等价的圈描述；负系数可用反向弧实现。

\ProseParagraphBreak
圈允许不同部分互相抵消。相同曲线正反各走一次就是零圈；大正圈减小正圈则是圆环的正向完整边界，内边在区域边界方向中为负向。把“每条边界都正向”理解成每条都逆时针，会在多连通区域里把内边符号写错。

\ProseParagraphBreak
圈是有限的整数线性组合，而不是一条必须首尾相接的参数路径。两个彼此不相交圆周可以组成一个圈，其积分、长度上界和指数都按系数相加。把它们用一条来回连接线串成单一闭路不会改变积分和指数，因为连接线的两个方向抵消；但这只是为画图或同伦描述方便，不是圈定义的一部分。有限性则保证这些求和与曲线像的紧性可直接使用。

\ProseParagraphBreak
系数取整数是为了保留指数的拓扑含义。若只做积分恒等式，当然可以考虑复系数线性组合；然而其“指数”不再是整数权重，不能直接解释成面片覆盖次数。一般留数公式可按线性延伸处理复系数，但零同调的拓扑定义与整数圈的几何填充应先建立，再作这样的代数扩展。

\subsection{同调}
\label{b3:subsec:100302}

\begin{definition}\label{b3:def:10-planar-homology}
区域 \(U\) 中圈 \(\Gamma_1,\Gamma_2\) 称为\term[tongdiao]{同调}[homologous]，若对每个 \(a\notin U\) 有 \(\operatorname{Ind}(\Gamma_1-\Gamma_2,a)=0\)。与零圈同调的圈称为在 \(U\) 中零同调。
\end{definition}
这是平面中一维整数同调的指数判别形式。下面的格网论证也说明：零同调圈在小路径替换后确可写为有限个包含于 \(U\) 的方形边界之和。因此定义具有边界相消的几何意义，不是把积分结论先作为定义。

\ProseParagraphBreak
同伦的闭路必同调，因为区域外点不在同伦像中，指数保持；反向一般不成立。积分只依靠同调的原因是每个局部方形边界积分为零，而方形相消不保留绕各孔的先后顺序。

\ProseParagraphBreak
同调忘掉闭路的非交换次序。在双孔域中，取基于同一点的绕孔闭路 \(\alpha,\beta\)。交换子 \(\alpha*\beta*\alpha^{-1}*\beta^{-1}\) 对每个外部点的指数为0，因为四段的指数成对抵消。它因此是零同调圈，任意全纯函数积分为0。但它在双孔域中不能收缩：将区域沿割线收缩到两个圈组成的图，基本群的约化词 \(\alpha\beta\alpha^{-1}\beta^{-1}\) 不消失。图的约化词判别是这里使用的标准拓扑事实。这个例子说明同伦形式虽容易表述，仍没有覆盖一般 Cauchy 零积分的全部情况。

\ProseParagraphBreak
圈的“填充”也不能只看图中的有界区域。设 \(U=\mathbb C\setminus\{0\}\)，两个同心正向圆周的差圈 \(\Gamma=\gamma_R-\gamma_r\) 对0的指数为0，故在 \(U\) 零同调；尽管每个圆单独都围住缺失点，两者合在一起只填充中间圆环。整数面片的系数恰记录这种相消，细网格证明中的 \(n_Q\) 就是各方格内的指数。负系数是代数定向，没有对应一个“负面积的区域”。

\subsection{零同调圈}
\label{b3:subsec:100303}

\begin{lemma}[格网分割]\label{b3:lem:10-grid-cycle}
若 \(\Gamma\) 在 \(U\) 中零同调，则它可在不改变任何全纯函数积分的条件下替换成一个细格网圈，并写成有限个闭方形 \(Q\subset U\) 的整数倍边界之和。
\end{lemma}
\begin{proof}
圈的支集紧含于 \(U\)，取小于它到补集距离的格网尺度。把各路径分成位于小圆盘内的段，先以折线再以邻近格边路径替换，接点用相同短连接弧对接。所有替换都在 \(U\) 的局部圆盘内同伦，局部原函数使积分不变，区域外指数也不变。所得格圈记为 \(G\)。

对每个开格面 \(Q\)，令 \(n_Q\) 为该面上的 \(\operatorname{Ind}(G,\cdot)\)。指数在面内恒定，无界远处为0，因此非零面有限。一条有向格边上的圈系数等于其左面指数减右面指数：在边中点附近横穿一次，连续辐角的总变化正跳过该条边的有向次数；也可对小横向两点的 \(1/(z-a)\) 积分逐边直接核对。因此按格边系数逐项比较，有 \(G=\sum_Q n_Q\partial Q\)。

若某个非零面的闭包含 \(a\notin U\)，则 \(a\) 不在格圈支集上。\(a\) 的小邻域指数恒为0，因为区域外指数为0；从该邻域进入相邻格面，迫使 \(n_Q=0\)，矛盾。故每个非零闭面包含于 \(U\)，完成分割。
\end{proof}
\begin{theorem}[同调形式 Cauchy 定理]\label{b3:thm:10-cauchy-homology}
若 \(f\) 在 \(U\) 全纯，\(\Gamma\) 在 \(U\) 中零同调，则 \(\int_\Gamma f\,\mathrm{d}z=0\)。对 \(a\in U\setminus\Gamma\)，还有
\[
 \frac1{2\pi\mathrm{i}}\int_\Gamma\frac{f(z)}{z-a}\,\mathrm{d}z
 =\operatorname{Ind}(\Gamma,a)f(a).
\]
\end{theorem}
\begin{proof}
格网分割将第一积分变成有限个小方形边界积分。每个闭方形在 \(U\) 内，其凸邻域上有原函数，故积分为0。第二式把被积式写成 \(f(a)/(z-a)+(f(z)-f(a))/(z-a)\)；第二项在 \(a\) 可去且全域全纯，由第一式积分为0，第一项按指数定义给所需值。
\end{proof}

\begin{corollary}[高阶同调积分公式]\label{b3:cor:10-higher-cauchy}
设 \(\Gamma\) 在 \(U\) 零同调，\(f\) 在 \(U\) 全纯，\(a\in U\setminus|\Gamma|\)，\(n\ge0\)。则
\[
 \frac1{2\pi\mathrm{i}}\int_\Gamma
 \frac{f(z)}{(z-a)^{n+1}}\,\mathrm{d}z
 =\operatorname{Ind}(\Gamma,a)\frac{f^{(n)}(a)}{n!}.
\]
\end{corollary}
\begin{proof}
\(n=0\) 是同调 Cauchy 公式。在 \(a\) 的一个不碰曲线的圆盘内，指数为常整数。积分核及其对参数的各阶导数在紧曲线与该小盘上连续有界，所以可在积分号下微分。对原公式连续微分 \(n\) 次，核导数为 \(n!/(z-a)^{n+1}\)，遂得结论。
\end{proof}
若 \(a\) 位于指数为0的区域，全部这些积分都为0，即使 \(f(a)\) 及其导数不为0。若圈绕 \(a\) 三次，所有结果乘3。局部估计从公式取得时，要同时记录曲线长度、曲线到 \(a\) 的距离及非零指数；不能对指数为0的圈把右边除以0来求 \(f(a)\)。

\ProseParagraphBreak
Cauchy 型积分还可以在曲线外定义一个函数：\(C(w)=(2\pi\mathrm{i})^{-1}\int_\Gamma f(z)/(z-w)\,\mathrm{d}z\)。它在每个补集分支全纯，在 \(U\setminus|\Gamma|\) 等于 \(\operatorname{Ind}(\Gamma,w)f(w)\)。无界分支的指数为0，故在那里为0。跨曲线时两侧整数权重可能不同，不能因此宣称 \(C\) 在整平面全纯；其定义本来就避开曲线。这个观察将积分表示与边界跃迁区分开来。

\ProseParagraphBreak
网格填充中的两个紧性条件值得分开。曲线像必须与 \(\mathbb C\setminus U\) 有正距离，才能把曲线逐段换成细网格边而不离开定义域；非零指数面片的闭包也必须留在 \(U\)，才能对每个方格使用 Cauchy 定理。后一个条件由圈在 \(U\) 零同调及指数局部常性保证，不是前一个条件的自动后果。例如围绕去心点的圆周与缺失点距离为正，但它的填充含该点，因而不能把整张填充都当成 \(U\) 内面片。

\ProseParagraphBreak
曲线可以通过分段光滑假设出现无限多次轻微振荡，仍不能简单按每次穿网格的时刻切分，交点可能无限。本章用紧路径的有限小盘链和分段近似替换来取得有限多边形圈，替换误差圈全部落在有局部原函数的小盘内，积分与外部指数都不变。再对有限多边形应用网格分解，面片个数有限。这种顺序避免把无限网格交叉当成有限组合来求和。

\subsection{拓扑留数定理}
\label{b3:subsec:100304}

\begin{theorem}[一般留数公式]\label{b3:thm:10-residue-general}
设 \(f\) 在区域 \(U\) 亚纯，\(\Gamma\) 在 \(U\) 中零同调且不经过 \(f\) 的极点。则
\[
 \frac1{2\pi\mathrm{i}}\int_\Gamma f(z)\,\mathrm{d}z
 =\sum_{a\in U}\operatorname{Ind}(\Gamma,a)\operatorname{Res}(f,a),
\]
其中只有有限个非零项。
\end{theorem}
\begin{proof}
指数非零的有界面及其闭包组成 \(U\) 内紧集：若闭包碰到区域外点，该点附近远离圈，指数局部恒为外部的0，矛盾。离散极点集在这个紧集只有有限点 \(a_j\)。各取互不交、避开圈的小正圆 \(c_j\)，令 \(m_j=\operatorname{Ind}(\Gamma,a_j)\)。圈 \(\Gamma-\sum_jm_jc_j\) 对 \(U\) 外点及所有极点的指数都为0，故在删去极点后的区域零同调。Cauchy 定理使其积分为0，每个 \(c_j\) 上的 Laurent 系数公式给 \(\int_{c_j}f=2\pi\mathrm{i}\operatorname{Res}(f,a_j)\)，代回即得结论。
\end{proof}
同一证明允许被积函数在这些孤立点为本性奇点，只需有关非零指数点有限，并在其余区域全纯。公式对极点的限制是常见且足够的亚纯版本。

\ProseParagraphBreak
例如圆环 \(r<|z|<R\) 的完整边界在稍大圆环中零同调，内外两圈对中心的指数相消；对全纯函数积分为0。但单独外圈不零同调，\(1/z\) 的积分非零。不能因圈“看起来是边界”便省去检查它是否是当前全纯区域中的边界。

\begin{corollary}[一般围道的辐角原理]\label{b3:cor:10-argument-general}
设 \(f\) 在 \(U\) 亚纯且不恒为0，\(\Gamma\) 在 \(U\) 零同调并避开 \(f\) 的零点与极点。则
\[
 \operatorname{Ind}(f\circ\Gamma,0)=
 \sum_a\operatorname{Ind}(\Gamma,a)\operatorname{ord}_a f.
\]
左边对圈逐段按整数系数理解，右边只有有限个非零项。
\end{corollary}
\begin{proof}
\(f'/f\) 在零极点处的留数为 \(\operatorname{ord}_a f\)。一般留数定理给右侧等于 \((2\pi\mathrm{i})^{-1}\int_\Gamma f'/f\,\mathrm{d}z\)；分段链式法则把它变成像圈的指数。零极点在相关的非零指数紧填充中只有有限多个，所以和有限。
\end{proof}
如果圈含负系数，式子是带权零极点计数。比如两个嵌套圆的差只计数圆环中的零极点，内孔的权重为0。这比只说“圈内的零点数减极点数”精确，因为一般圈没有单一的内部。

\ProseParagraphBreak
还可以在核前乘一个全纯函数 \(h\)，得到加权零极点和
\[
 \frac1{2\pi\mathrm{i}}\int_\Gamma h(z)\cdot\frac{f'(z)}{f(z)}\,\mathrm{d}z
 =\sum_a\operatorname{Ind}(\Gamma,a)\cdot h(a)\cdot\operatorname{ord}_a f.
\]
选 \(h(z)=z^k\) 就计算带重数零极点的幂和。若只有一个指数为1的 Jordan 围道且里面无极点，\(k=1\) 给零点总和，\(k=0\) 给零点总数。计算需避开围道零点，不是对含零点的边界积分取未声明的主值。

\ProseParagraphBreak
割线围道中的两岸方向尤其容易出错。对正实轴割开的平面取 \(0<\arg z<2\pi\)，上岸由内半径走向外半径，下岸由外走向内；两岸的 \(z^{s-1}\) 因子相差 \(e^{2\pi\mathrm{i}(s-1)}\)。因此围道积分中两岸合成为 \((1-e^{2\pi\mathrm{i}s})\int_r^R x^{s-1}/(1+x)\,\mathrm{d}x\)，而不是两项直接相加。小圆弧的估计为 \(O(r^{\operatorname{Re}s})\)，大圆弧为 \(O(R^{\operatorname{Re}s-1})\)，恰需 \(0<\operatorname{Re}s<1\)。位于 \(-1\) 的留数按分支取 \(e^{\pi\mathrm{i}(s-1)}\)。整理得
\[
 \int_0^\infty\frac{x^{s-1}}{1+x}\,\mathrm{d}x=\frac\pi{\sin\pi s}.
\]
这个计算先在两岸略离割线的有限围道使用留数定理，再令岸逼近、半径趋于0和无穷；各项估计说明极限合法。后面特殊函数中会再次出现此积分，届时使用的是已确认的分支与收敛范围。

\ProseParagraphBreak
一般留数定理并不要求整个 \(U\) 只有有限个极点。相关圈的非零指数区域被一个紧子集包住，亚纯函数的极点在紧子集中只有有限个；圈外的其他极点权重为0，不出现在和中。若极点向 \(U\) 的边界聚集，仍与局部亚纯性相容。若它们在 \(U\) 内聚集，就不能是通常意义的亚纯函数，公式的假设已失败。

\ProseParagraphBreak
留数在小圈上的计算与大圈上的拓扑权重分工明确。局部 Laurent 系数决定每个奇点贡献多少，指数决定围道取该贡献几次。先局部计算再乘权重，可处理嵌套圈、负向圈和多重绕圈；如果从图像直接判断一个奇点“在里面”，容易漏掉权重0、负权重或大于1的情形。计算过程还应保留奇点恰落在曲线上的例外，此时通常积分未定义，需要改变围道或另行声明主值问题。

\ProseParagraphBreak
拓扑条件也不能用“所有奇点都有指数0”在任意函数上替代。对一个特定亚纯函数，某些缺失点未必是它的奇点；即使其已知奇点的加权留数为0，函数也可能在未填入的定义域外没有任何指定延拓。一般 Cauchy 定理的零同调要求针对整个 \(\mathbb C\setminus U\)，保证填充合法；针对某一函数的特殊抵消可以给额外零积分，但应通过该函数的实际延拓或直接计算证明。

\ProseParagraphBreak
本章从路径同伦走向整数圈同调，没有引入高维同调的全部定义。这里的指数判别与平面面片填充已经足够证明需要的 Cauchy 公式。一般流形上的链、边界算子与同调群会把这个过程抽象化；若在曲面上没有全局平面坐标，就不能继续以“外部点的绕数”直接定义同调，必须改用那套链理论。附录A的曲面留数证明通过局部三角剖分相消，正是保持同一积分思想并换成适合曲面的对象。
